\begingroup
\setlength{\parindent}{0pt}
\setlength{\parskip}{1.0\baselineskip}
\setlength{\abovedisplayskip}{8pt}
\setlength{\belowdisplayskip}{8pt}
\setlength{\abovedisplayshortskip}{8pt}
\setlength{\belowdisplayshortskip}{8pt}
\setlength{\jot}{5pt}
\clubpenalty=10000
\widowpenalty=10000
\displaywidowpenalty=10000
\makeatletter
\let\NSNativeDisplay\[
\long\def\[#1\]{%
  \par\nobreak\vskip8pt\nointerlineskip
  \begingroup
  \parskip=0pt \predisplaypenalty=10000 \postdisplaypenalty=10000
  \baselineskip=0pt \lineskip=0pt \lineskiplimit=0pt
  \abovedisplayskip=0pt \belowdisplayskip=0pt
  \abovedisplayshortskip=0pt \belowdisplayshortskip=0pt
  \NSNativeDisplay \relax{\normalbaselines #1}\]
  \par
  \endgroup
  \nobreak\vskip8pt\nointerlineskip
  \parskip=0pt
  \AddToHookNext{para/begin}{\ifdefined\NSNativeDisplay\setlength{\parskip}{1.0\baselineskip}\fi}%
  \ignorespaces
}
\makeatother

\section*{How the Proof Works}
\label{guide:how-proof-works}

The three-dimensional Navier--Stokes regularity problem asks whether a smooth fluid flow can lose its smoothness in finite time. We start with smooth, divergence-free initial velocity \(u_0\), finite energy, and positive viscosity. The fluid evolves according to
\[
\partial_tu+(u\cdot\nabla)u=-\nabla p+\nu\Delta u,
\qquad \nabla\cdot u=0.
\]

The problem starts with a vortex.

Strain within the surrounding flow can stretch a vortex. As the vortex lengthens, incompressibility makes it thinner, and the fluid can spin faster as its rotation concentrates into the smaller cross-section. Viscosity spreads the vorticity while the surrounding flow continues pulling on the vortex. Even if the vortex loses circulation, its rotation can intensify if its cross-section shrinks fast enough. The narrowing and the viscous spreading happen together, and the danger lies in how their rates develop.

Physically, we never witness a singularity in a real fluid. What we actually
see is a vortex spinning faster and faster, getting thinner and thinner as
it stretches, until it dissipates into the surrounding fluid and disappears
altogether.

So the intuition is that, perhaps, as the vortex thins out and speeds up,
viscosity does something that allows it to overcome the nonlinear transport.
But the equations describe an idealized continuum, with no minimum radius
built into the model. We have to allow for the possibility that the vortex
keeps narrowing toward an infinitesimally small radius. To study what happens
there, we rescale the equations and examine how viscosity and nonlinear
transport behave as the vortex becomes smaller and faster.

Contract lengths by a factor \(\lambda>1\), increase the velocity by
\(\lambda\), and run the motion \(\lambda^2\) times as fast:
\[
 u_\lambda(x,t)=\lambda u(\lambda x,\lambda^2t),
 \qquad p_\lambda(x,t)=\lambda^2p(\lambda x,\lambda^2t).
\]
All terms in the momentum equation scale by \(\lambda^3\), preserving
the balance between transport and viscosity at fixed \(\nu\). The velocity
gradients scale by \(\lambda^2\), since the increase in speed is compressed
into a smaller distance.

Yet the energy decreases, because the \(\lambda^2\) increase in velocity
squared is outweighed by the \(\lambda^{-3}\) decrease in volume:
\[
 E[u_\lambda](t)=\lambda^{-1}E[u](\lambda^2t),
 \qquad E[u]=\frac12\int_{\mathbb R^3}|u|^2\,dx.
\]

A Sobolev norm includes the growing spatial derivatives in the measurement.
Its order \(s\) contributes \(\lambda^s\), alongside the
\(\lambda^{-1/2}\) from the velocity's \(L^2\) scaling:
\[
 \|u_\lambda(\cdot,t)\|_{\dot H^s}
 =\lambda^{s-1/2}\|u(\cdot,\lambda^2t)\|_{\dot H^s}.
\]
Below \(s=\tfrac12\), the norm decreases under contraction; above it,
the norm increases. At the critical height \(s=\tfrac12\), the growing
velocity and spatial variation exactly compensate for the shrinking
volume, leaving the norm unchanged.

If each contraction generated another, the time rescaling would give the
stages progressively shorter durations:
\[
 \tau_j=\lambda^{-2j}\tau_0,
 \qquad \sum_{j=0}^{\infty}\tau_j
 =\frac{\tau_0}{1-\lambda^{-2}}<\infty.
\]
The infinitely many stages would accumulate at a finite time, with the
radius tending to zero and the velocity and its gradients growing without
bound. This is the hypothetical Zeno cascade, and the critical height now
looks like finite-time blowup.

If a singularity were to actually form, the velocity would have to blow up,
which means the acceleration would have to blow up, followed by its jerk,
and so on, and so on.

Differentiate the equation in time at fixed spatial coordinates, once
and then again. Each time derivative of transport can fall on either
velocity factor, giving
\[
 \begin{aligned}
 &\partial_tu=\nu\Delta u-(u\cdot\nabla)u-\nabla p,\\
 &\partial_t^2u=\nu\Delta\partial_tu-(\partial_tu\cdot\nabla)u
   -(u\cdot\nabla)\partial_tu-\nabla\partial_tp,\\
 &\partial_t^3u=\nu\Delta\partial_t^2u-(\partial_t^2u\cdot\nabla)u
   -2(\partial_tu\cdot\nabla)\partial_tu
   -(u\cdot\nabla)\partial_t^2u-\nabla\partial_t^2p.
 \end{aligned}
\]
Continuing gives the temporal derivative tower. Writing the derivatives at
fixed spatial coordinates as \(V_n=\partial_t^nu\) and
\(Q_n=\partial_t^np\), its general row is
\[
 \begin{aligned}
 V_{n+1}&=\nu\Delta V_n
 -\sum_{a=0}^{n}\binom na(V_a\cdot\nabla)V_{n-a}-\nabla Q_n,\\
 -\Delta Q_n&=\partial_i\partial_j
   \sum_{a=0}^{n}\binom na V_{a,i}V_{n-a,j}.
 \end{aligned}
\]
Taking the divergence fixes the pressure equation, and decay at spatial
infinity fixes its normalization. Every temporal row is also a spatial
field. Taking its spatial derivatives gives the mixed jet
\(J_{n,\alpha}=\partial_x^\alpha\partial_t^nu\), with pressure derivatives
\(\Pi_{n,\alpha}=\partial_x^\alpha\partial_t^np\). The product rule now
acts in both directions:
\[
 J_{n+1,\alpha}=\nu\Delta J_{n,\alpha}-\nabla\Pi_{n,\alpha}
 -\sum_{a=0}^{n}\sum_{\beta\le\alpha}
   \binom na\binom\alpha\beta
   (J_{a,\beta}\cdot\nabla)J_{n-a,\alpha-\beta}.
\]
The mixed jet keeps the velocity and pressure derivatives as fields.
We now take a scalar energy measurement of \(u\):
\[
 E_s=\frac12\|\Lambda^su\|_2^2,
 \qquad \Lambda=(-\Delta)^{1/2}.
\]
The row \(n=0,\alpha=0\) is the Navier--Stokes equation for \(u\).
Differentiating \(E_s\), substituting that row, and integrating by parts gives
\[
 \begin{aligned}
 E_s'
 &=\operatorname{Re}\langle\Lambda^su,\Lambda^s\partial_tu\rangle\\
 &=-\nu\|\Lambda^{s+1}u\|_2^2+\mathcal P_s\\
 &=-2\nu E_{s+1}+\mathcal P_s,
 \end{aligned}
\]
\Needspace{5\baselineskip}
where the transport and pressure contributions are
\[
 \mathcal P_s=-\operatorname{Re}\langle\Lambda^su,
 \Lambda^s((u\cdot\nabla)u+\nabla p)\rangle.
\]
One time derivative of the energy has brought in an energy one spatial
Sobolev level higher. The Laplacian contributes two spatial derivatives;
integration by parts places one additional derivative on each velocity
factor in the squared norm.

At the critical height, set \(H=E_{1/2}\). Repeating the energy identity
at each newly exposed height gives
\[
 \begin{aligned}
 H'&=-2\nu E_{3/2}+\mathcal P_{1/2},\\
 H''&=-2\nu E_{3/2}'+\partial_t\mathcal P_{1/2}\\
    &=4\nu^2E_{5/2}-2\nu\mathcal P_{3/2}
        +\partial_t\mathcal P_{1/2},\\
 H^{(n)}&=(-2\nu)^nE_{n+1/2}
        +\sum_{j=0}^{n-1}(-2\nu)^j
           \partial_t^{\,n-1-j}\mathcal P_{j+1/2}.
 \end{aligned}
\]
The higher mixed-jet rows organize the velocity and pressure derivatives
inside \(\partial_t^k\mathcal P_s\). Here \(H^{(n)}\) differentiates the
critical energy; it is not the energy of \(\partial_t^nu\).

Each step up in the time derivatives of \(H\) brings in one higher spatial
Sobolev order. At temporal order \(n\), the energy identity exposes
\(E_{n+1/2}\), so the sequence \(H,H',H'',\ldots\) takes us through
the spatial levels
\[
 H^{1/2},\qquad H^{3/2},\qquad H^{5/2},\qquad H^{7/2},\qquad\ldots.
\]

\phantomsection\label{guide:intersection}
Recall the Sobolev embedding theorem: in three dimensions, enough spatial
Sobolev regularity gives bounded, continuous derivatives through order \(k\):
\[
 H^s(\mathbb R^3)\hookrightarrow C_b^k(\mathbb R^3),
 \qquad s>k+\tfrac32.
\]

As we go up the temporal tower, membership at the corresponding spatial
levels gives more regularity and brings us closer to smoothness. A vortex
approaching blowup would make us expect the opposite: its spatial derivatives
should be growing without bound. Requiring membership at every height gives
the full smoothness condition:
\[
 L^2\cap\bigcap_{n\ge0}\dot H^{n+1/2}(\mathbb R^3)
 =\bigcap_{m\ge0}H^m(\mathbb R^3)
 =H^\infty(\mathbb R^3)\hookrightarrow C_b^\infty(\mathbb R^3).
\]

\Cref{thm:compatible-inverse-intersection} establishes this simultaneous
regularity in space and time for the flow from
\(u_0\in H^\infty_\sigma(\mathbb R^3)\), throughout the interval leading
to a proposed finite maximal time \(T_*\):
\[
 u\in\bigcap_{r,m\ge0}H^r(0,T_*;H^m(\mathbb R^3)).
\]

At finite orders, this intersection gives a rectangle of the mixed jet.
Choose a temporal depth \(N\) and spatial height \(M\). The rectangle
contains the rows \(V_0,\ldots,V_N\), their next time derivatives, and
pressure:
\[
 V_n\in L^2(0,T_*;H^{M+1}),\qquad
 V_{n+1}\in L^2(0,T_*;H^{M-1}),\qquad 0\le n\le N,
\]
with \(Q_n\in L^1(0,T_*;H^M)\).

\Needspace{8\baselineskip}
Viscosity gives the relationship between these rows. The term
\(\nu\Delta V_n\) connects the next time derivative to two more spatial
derivatives, with transport and pressure retained. Positive \(\nu\) lets us
recover that spatial contribution from each complete row. The spreading we
pictured at the start is viscosity acting on differences in velocity across
the vortex; in the energy identity, its contribution
\(-\nu\|\Lambda^{s+1}u\|_2^2\) measures the spatial variation that is
becoming steeper as the vortex narrows.

In particular, choose a height large enough to put \(V_{n+1}\) in
\(L^2(0,T_*;H^m)\). Between two late times,
\[
 \|V_n(t)-V_n(s)\|_{H^m}
 \le(t-s)^{1/2}\|V_{n+1}\|_{L^2(s,t;H^m)},
 \qquad s<t<T_*.
\]
The remaining change tends to zero, so the row reaches an \(H^m\) limit.
Limits at different spatial heights agree under restriction. For the
velocity they identify a smooth endpoint
\(u_*\in H^\infty_\sigma(\mathbb R^3)\), with the higher temporal rows
and normalized pressure reaching their compatible endpoints. Local existence
continues from \(u_*\), and uniqueness joins that continuation to the
incoming flow. A finite \(T_*\) cannot be maximal.

\phantomsection\label{guide:finite-estimates}
The finite estimates now come from whichever rectangle contains the
orders they require. For temporal order \(r\), spatial derivative order
\(k\), and finite time \(T\), choose \(m>k+\tfrac32\). Continuity in
\(H^m\) and Sobolev embedding give
\[
 \sup_{0\le t\le T}\|\partial_t^ru(t)\|_{W^{k,\infty}}
 \le C_{m,k}\sup_{0\le t\le T}\|\partial_t^ru(t)\|_{H^m}<\infty.
\]
In particular, the accumulated velocity gradient is finite. For two material
trajectories, this gives
\[
 |X(a,t)-X(b,t)|\ge |a-b|
 \exp\!\left(-\int_0^t\|\nabla u(s)\|_\infty\,ds\right)>0.
\]
The vortex can stretch and narrow, but a positive material separation
retains a positive lower bound throughout each finite interval. The
intersection supplies the regularity, and its finite rectangles supply
the estimates at the orders where we need them.

So, for example, choose the rectangle with \(N=0\) and \(M=1\).
Write the sizes of its velocity and time-derivative entries as
\[
 \begin{aligned}
 A_T&=\left(\int_0^T\|u(t)\|_{H^2}^2\,dt\right)^{1/2},\\
 B_T&=\left(\int_0^T\|\partial_tu(t)\|_2^2\,dt\right)^{1/2}.
 \end{aligned}
\]
The time-derivative entry controls the change in the \(H^1\) trace.
Pairing \(\partial_tu\) with \((1-\Delta)u\), then applying
Cauchy--Schwarz in time, gives, for \(0\le t\le T\),
\[
 \begin{aligned}
 \|u(t)\|_{H^1}^2
 &=\|u_0\|_{H^1}^2
   +2\operatorname{Re}\int_0^t
       \langle(1-\Delta)u(s),\partial_su(s)\rangle_{L^2}\,ds\\
 &\le\|u_0\|_{H^1}^2+2A_TB_T.
 \end{aligned}
\]
\Needspace{6\baselineskip}
With the unitary Fourier normalization, the quantitative spatial
\(H^1\hookrightarrow L^3\) estimate is
\[
 \|f\|_{L^3}\le\frac{\|f\|_{H^1}}{2^{5/6}\pi^{1/6}}.
\]
Applying it to the trace just obtained gives the
Escauriaza--Seregin--\v{S}ver\'ak endpoint estimate:
\[
 \sup_{0\le t\le T}\|u(t)\|_{L^3}
 \le \frac{\bigl(\|u_0\|_{H^1}^2+2A_TB_T\bigr)^{1/2}}
              {2^{5/6}\pi^{1/6}}.
\]

The \(H^2\) entry also controls peak speed at each time. The spatial
estimate is
\[
 \|u(t)\|_{L^\infty}\le\frac{\|u(t)\|_{H^2}}{\sqrt{8\pi}}.
\]
Squaring and integrating gives the Serrin estimate:
\[
 \left(\int_0^T\|u(t)\|_{L^\infty}^2\,dt\right)^{1/2}
 \le\frac1{\sqrt{8\pi}}
       \left(\int_0^T\|u(t)\|_{H^2}^2\,dt\right)^{1/2}
 =\frac{A_T}{\sqrt{8\pi}}.
\]

For the gradient, the spatial \(H^2\hookrightarrow W^{1,3}\) estimate gives
\[
 \|\nabla u(t)\|_{L^3}
 \le\frac{\sqrt3\,7^{1/6}}{2^{5/3}\pi^{1/6}}\|u(t)\|_{H^2}.
\]
Integrating its square gives the Beir\~ao da Veiga estimate:
\[
 \left(\int_0^T\|\nabla u(t)\|_{L^3}^2\,dt\right)^{1/2}
 \le\frac{\sqrt3\,7^{1/6}}{2^{5/3}\pi^{1/6}}
       \left(\int_0^T\|u(t)\|_{H^2}^2\,dt\right)^{1/2}
 =\frac{\sqrt3\,7^{1/6}}{2^{5/3}\pi^{1/6}}A_T.
\]
\phantomsection\label{guide:how-proof-ends}
The rectangle supplies all three estimates. Choosing higher rows and
spatial heights extends the calculation to higher derivatives.

\endgroup
